\documentclass[12pt]{article}
\usepackage[utf-8]{inputenc}
\usepackage{graphicx}
\usepackage{hyperref}
\usepackage{listings}
\usepackage{xcolor}
\usepackage{geometry}

\geometry{margin=1in}

\lstset{
    language=Python,
    basicstyle=\ttfamily\small,
    keywordstyle=\color{blue},
    commentstyle=\color{gray},
    stringstyle=\color{red},
    breaklines=true,
    showstringspaces=false
}

\title{AutoShip: Automated Python Code Quality Analysis System}
\author{Development Team}
\date{\today}

\begin{document}

\maketitle

\tableofcontents

\newpage

\section{Introduction}

AutoShip is a web-based application designed to automate Python code quality analysis and reporting. The system leverages FastAPI for the backend and integrates with pylint for comprehensive code quality metrics.

\section{Project Overview}

\subsection{Objectives}

\begin{enumerate}
    \item Provide automated code quality analysis for Python projects
    \item Track analysis history and trends over time
    \item Generate detailed reports with actionable insights
    \item Enable team collaboration through shared dashboards
    \item Maintain comprehensive audit logs for compliance
\end{enumerate}

\subsection{Key Features}

\begin{itemize}
    \item Real-time code quality scoring
    \item Historical analysis tracking
    \item Pass/fail assessment
    \item Trend analysis and visualization
    \item Database persistence using SQLAlchemy ORM
    \item Responsive web interface with Jinja2 templating
    \item RESTful API endpoints
\end{itemize}

\section{Technology Stack}

\subsection{Backend}
\begin{itemize}
    \item \textbf{FastAPI}: Modern Python web framework for building APIs
    \item \textbf{SQLAlchemy}: SQL toolkit and Object-Relational Mapping (ORM)
    \item \textbf{Pylint}: Python code quality analysis tool
    \item \textbf{Uvicorn}: ASGI server for running FastAPI applications
\end{itemize}

\subsection{Frontend}
\begin{itemize}
    \item \textbf{HTML5}: Markup language for web pages
    \item \textbf{CSS3}: Styling and responsive design
    \item \textbf{Jinja2}: Template engine for dynamic content rendering
\end{itemize}

\subsection{Deployment}
\begin{itemize}
    \item \textbf{Docker}: Containerization for consistent environments
    \item \textbf{Docker Compose}: Multi-container orchestration
\end{itemize}

\section{System Architecture}

\subsection{Directory Structure}

\begin{lstlisting}
AutoShip/
├── app/
│   ├── main.py              # FastAPI application entry point
│   ├── models.py            # SQLAlchemy ORM models
│   ├── database.py          # Database configuration
│   ├── static/              # CSS and client-side assets
│   └── templates/           # Jinja2 HTML templates
├── Dockerfile               # Container configuration
├── docker-compose.yml       # Multi-container setup
├── requirements.txt         # Python dependencies
└── report.tex              # This documentation
\end{lstlisting}

\subsection{Core Components}

\subsubsection{Database Models}
The system uses SQLAlchemy ORM with the following primary model:
\begin{itemize}
    \item \textbf{AnalysisResult}: Stores code analysis results with timestamps, scores, and status flags
\end{itemize}

\subsubsection{API Endpoints}
\begin{itemize}
    \item \texttt{GET /}: Dashboard with statistics and recent analyses
    \item \texttt{GET /health}: Health check endpoint
    \item Analysis submission and retrieval endpoints
\end{itemize}

\subsubsection{Templates}
\begin{itemize}
    \item \texttt{base.html}: Base template with common layout
    \item \texttt{index.html}: Landing page and dashboard
    \item \texttt{analyze.html}: Code submission interface
    \item \texttt{results.html}: Detailed analysis results
    \item \texttt{history.html}: Historical analysis records
\end{itemize}

\section{Implementation Details}

\subsection{Data Flow}

\begin{enumerate}
    \item User submits Python code through the web interface
    \item FastAPI receives the submission and validates input
    \item Pylint analyzes the code and generates quality score
    \item Results are persisted to the database via SQLAlchemy ORM
    \item Analysis history is tracked with timestamps
    \item Results are rendered and displayed to the user
\end{enumerate}

\subsection{Key Statistics}

The dashboard calculates and displays:
\begin{itemize}
    \item Total number of analyses performed
    \item Average pylint score across all analyses
    \item Pass rate (percentage of analyses meeting quality threshold)
    \item Today's analysis count
    \item Recent analyses (last 5)
    \item Score trend data (last 10 analyses)
\end{itemize}

\section{Development and Deployment}

\subsection{Local Development}

Requirements:
\begin{itemize}
    \item Python 3.8 or higher
    \item FastAPI and dependencies (see requirements.txt)
    \item A supported database (SQLite, PostgreSQL, etc.)
\end{itemize}

\subsection{Docker Deployment}

The project includes Docker configuration for containerized deployment:
\begin{itemize}
    \item \texttt{Dockerfile}: Defines the application container
    \item \texttt{docker-compose.yml}: Orchestrates services
\end{itemize}

Run with:
\begin{lstlisting}
docker-compose up --build
\end{lstlisting}

\section{Future Enhancements}

\begin{enumerate}
    \item Integration with additional code quality tools (flake8, black, mypy)
    \item Advanced filtering and search capabilities
    \item Export functionality (PDF, CSV, JSON)
    \item User authentication and role-based access control
    \item CI/CD pipeline integration
    \item Real-time notifications for quality regressions
    \item Performance optimization and caching layer
\end{enumerate}

\section{Conclusion}

AutoShip provides a comprehensive solution for automated Python code quality analysis. By combining modern Python frameworks with robust ORM practices, the system offers scalability, maintainability, and extensibility for team-based development workflows.

\end{document}
